% !TEX root = ../main.tex
% ---------------------------------------------------------------------
%  P2 -- RISC-V Vektorerweiterung
% ---------------------------------------------------------------------
\vlsec{cP2}{P2}{SIMD: RISC-V Vektorerweiterung}
\label{sec:rvv}

\begin{formel}[Der eine Unterschied zu NEON]
Die Vektorlänge steht \emph{nicht} im Befehlssatz. Derselbe Binärcode
läuft auf 2-, 4- oder 8-Lane-Hardware -- er braucht nur mehr oder weniger
Takte. Keine Neukompilierung, kein Multiversioning.
\end{formel}

\begin{itemize}
\item 32 Vektorregister \R{v0}--\R{v31}; ein Vektor hält Skalar, Vektor
      oder Matrix
\item \T{VL-Register} definiert die benutzte Länge, $\text{VL}\le\text{MAXVL}$
\item \R{vfadd.s v0,v1,v2}: für $i<\text{VL}$ addieren,
      für $\text{VL}\le i<\text{MAXVL}$ Nullen schreiben
\item \R{vconfig} setzt Datenbreite und Registerzahl und schaltet die
      Einheit ab (Strom!)
\end{itemize}

\subsection{Strip Mining: Tail-Problem in Hardware gelöst}
\begin{lstlisting}
vconfig 0x63        # 32b Daten, 4 vregs, VL=MAXVL
vfld.s  v1, 0(x4)   # Konstante laden + broadcast
loop:
  vsetvl x2, x1     # VL = x2 = min(MAXVL, Rest x1)
  vflw   v0, 0(x3)  # laedt genau VL Elemente
  vfadd.s v2, v1, v0
  vsw    v2, 0(x3)
  slli x5,x2,2      # Bytes = VL * sizeof(float)
  add  x3,x3,x5     # Zeiger weiter
  sub  x1,x1,x2     # Rest verringern
  bne  x1,x0,loop
vconfig 0x0         # Vektoreinheit aus
\end{lstlisting}
\R{vsetvl} liefert $\min(\text{MAXVL},\text{Rest})$ \ra die letzte
Iteration verarbeitet automatisch nur den Rest. \T{Kein Tail-Code.}

\subsection{Befehlsklassen}
\begin{ztab}{@{}lX@{}}
Load/Store & \R{vlw/vsw} unit-stride \\
 & \R{vls*/vss*} \T{strided} (jedes 2., 3., \dots{} Element) \\
 & \R{vlx*/vsx*} \T{indexed} $=$ Gather/Scatter \\
Integer & \R{vadd}, \R{vmul}, \T{\R{vdiv}}, \R{vsll}, \R{vand},
  \R{vseq/vslt}, \R{vclip} (fixed point) \\
Float & \R{vfadd/sub/mul/div}, \R{vfsqrt}, \R{vfmax/min},
  \R{vfclass}; Suffix \R{.h/.s/.d} $=$ FP16/32/64 \\
MAC & \R{vfmadd}, \R{vfmsub}, widening \R{vfmaddwdn} \\
Convert & \R{vfcvt.<ziel>.<quelle>} int$\leftrightarrow$float,
  float$\leftrightarrow$float \\
Movement & \R{vins}/\R{vext} (GPR$\leftrightarrow$Lane),
  \R{vmerge} (maskengesteuert), \R{vrgather}
  (\R{vd[i]=vs1[vs2[i]]}), \R{vsplat}, \R{vslideup/dwn} \\
Masken & \R{vmfirst}, \R{vmpopc}, \R{vmiota} \\
\end{ztab}

\subsection{Gegenüberstellung}
\begin{ztab}{@{}lXX@{}}
 & NEON & RISC-V V \\ \midrule
Länge & fix 128\,Bit & VL variabel \\
Tail & Hand-Code & \R{vsetvl} \\
Masken & Behelf (\R{VBSL}) & eigene Maskenregister \\
Gather/Scatter & -- & \R{vlx*}/\R{vsx*} \\
Division & Newton-Raphson & \R{vdiv} \\
FP64 & -- & \R{.d} \\
\end{ztab}
Die RVV-Erweiterung behebt also genau die vier Schwachstellen, die bei
NEON in Abschnitt~\ref{sec:simdgrenzen} als Grenzen stehen.
